\documentclass[11pt]{article}
\usepackage[margin=1in]{geometry}
\usepackage{amsmath,amssymb,amsthm,booktabs}
\usepackage[hidelinks]{hyperref}
\newtheorem{theorem}{Theorem}
\newtheorem{lemma}{Lemma}
\newcommand{\D}{\mathbb D}
\newcommand{\ph}{\mathsf p}
\title{A covering-strip and moment certificate for Korenblum's constant}
\author{Ricardo Pinho\\\small ORCID: \href{https://orcid.org/0009-0006-2744-2599}{0009-0006-2744-2599}}
\date{Unrefereed preprint, 19 October 2026}
\begin{document}
\maketitle
\begin{abstract}
We give an explicit certificate for $c_2\ge54/125=0.432$, where $c_2$ is
Korenblum's constant for the Bergman space. This is an unrefereed preprint. The construction combines an
elementary two-zero Schwarz estimate on the covering strip of an annulus,
the classical hyperbolic Jung inequality, and Wikstr\"om's moment-duality
criterion. A rational twelve-atom measure with a Lebesgue tail satisfies the
required inequalities for every moment index. Two independently implemented
exact checkers, one using Arb and one using only integer interval arithmetic,
verify the quantitative hypotheses. The best public certified lower bound
located in the accompanying priority audit is Wikstr\"om's $0.4263$.
Neither optimality nor a new version of Jung's theorem is claimed.
\end{abstract}

\section{Statement and analytic reduction}
Korenblum introduced the maximum principle in~\cite{Korenblum};
Hayman~\cite{Hayman} proved that a positive admissible radius exists.
We give a lower bound, not a determination of the optimal constant.

Normalize Bergman area by $dA=dx\,dy/\pi$, so
$\|\sum a_kz^k\|_2^2=\sum |a_k|^2/(k+1)$. Korenblum's constant $c_2$ is
the supremum of those $c\in(0,1)$ for which $|f|\le|g|$ on
$c<|z|<1$ implies $\|f\|_2\le\|g\|_2$, for all $f,g\in A^2(\D)$.

\begin{theorem}\label{thm:main}
Korenblum's constant satisfies $c_2\ge54/125$.
\end{theorem}

We retain the moment framework of Wikstr\"om~\cite{Wikstrom}, but use a
separately derived coefficient estimate. Put $w=-\log c$,
\[
 X=\sinh^2\frac{\pi^2}{2w},\qquad
 K(t)=\frac{X}{X+\sin^2\!\left(\frac{\pi\log t}{2w}\right)},\qquad
 G(t)=\frac12\left(1+\sqrt{\frac{3+K(t)^2}{3(1-K(t)^2)}}\right)
 \quad(c^2<t<1).
\]
Write $\beta(t)=1-1/G(t)$, with $\beta(c^2)=\beta(1)=1$ by limits.
The $K$ here is a covering-strip bound; it is not the annular product
appearing in Wang's estimate as quoted in~\cite{Wikstrom}.

\begin{lemma}\label{lem:diam}
If $\omega:A(c,1)\to\D$ is holomorphic, then its image of $|z|=r$
has pseudohyperbolic diameter at most $K(r^2)$.
\end{lemma}
\begin{proof}
Lift to $h(\zeta)=\omega(e^\zeta)$ in $\log c<\operatorname{Re}\zeta<0$.
The map $\exp(i\pi(\zeta-\log c)/w)$ maps this strip biholomorphically
onto the upper half-plane. After a rotation of the annulus, choose lifts
$a=\log r$ and $a+i\theta$, $0\le\theta\le2\pi$, of the two points.
Postcompose $h$ with a disk automorphism making its value at $a$ zero.
Periodicity gives another zero at $a+2\pi i$. Dividing successively by
the two Blaschke factors and using Schwarz's lemma bounds the value at the
other point by the product of its pseudohyperbolic distances to these zeros.

Set $A=\pi^2/w$, $\phi=\pi\log(r/c)/w$ and $v=\pi\theta/w$.
The distances in question are $\ph(v)$ and $\ph(2A-v)$, where
\[
 \ph(v)=\frac{\sinh(v/2)}{\sqrt{\sinh^2(v/2)+\sin^2\phi}}.
\]
For $u=v/2>0$, the derivative of $\log\ph(2u)$ is
\[
 \frac{\sin^2\phi\,\coth u}{\sinh^2u+\sin^2\phi}.
\]
This is positive and decreasing, because its positive numerator decreases
and its positive denominator increases. Thus $\log\ph$ is concave, and
$\ph(v)\ph(2A-v)\le\ph(A)^2=K(r^2)$. Endpoints follow by continuity.
\end{proof}

\begin{lemma}\label{lem:coeff}
Suppose $f=\sum a_kz^k$, $g=\sum b_kz^k$ are holomorphic in $\D$ and
$|f|\le|g|$ on $A(c,1)$. For $c<r<1$, let
\[
 Y(r^2)=\sum_{\substack{k\ge1\\|a_k|<|b_k|}}
 (|b_k|^2-|a_k|^2)r^{2k},\qquad
 M_h(r)^2=\frac1{2\pi}\int_0^{2\pi}|h(re^{i\theta})|^2\,d\theta.
\]
Then $Y(r^2)\le G(r^2)(M_g(r)^2-M_f(r)^2)$.
\end{lemma}
\begin{proof}
The quotient $f/g$ has removable singularities in the annulus. The constant
and boundary cases are immediate, so assume it maps into $\D$.
Equip $\D$ with hyperbolic distance twice the inverse hyperbolic tangent of
the pseudohyperbolic distance. By Lemma~\ref{lem:diam}, the compact image of
the circle has hyperbolic diameter at most $D=2\operatorname{artanh}K(r^2)$.
The classical hyperbolic Jung inequality~\cite{BMRS} places it in a disk of
hyperbolic radius $R_h$ with
\[
 \sinh R_h\le\frac2{\sqrt3}\sinh(D/2).
\]
For completeness, consider a smallest enclosing hyperbolic disk. At its
center the directions to its boundary contacts cannot lie in an open
semicircle, since a small motion would shrink the radius. Two opposite
contacts, or three contacts whose directions contain zero in their convex
hull, give a pair separated by an angle at least $2\pi/3$. The hyperbolic
cosine rule gives $\cosh D\ge1+\tfrac32\sinh^2R_h$, proving the inequality.

The enclosing disk has pseudohyperbolic center $\alpha\in\D$ and radius
$\rho=\tanh(R_h/2)$, and hence
\[
 \frac1{1-\rho^2}=\frac{1+\cosh R_h}{2}\le G(r^2).
\]
Set $F=f-\alpha g$ and $H=g-\overline\alpha f$. On the circle,
$|F|\le\rho|H|$, and coefficientwise
\[
 |H_k|^2-|F_k|^2=(1-|\alpha|^2)(|b_k|^2-|a_k|^2).
\]
Dropping the negative $|F_k|^2$ contributions in the indices defining $Y$
and adding nonnegative $|H_k|^2$ terms gives
\[
 Y(r^2)\le\frac{M_H(r)^2}{1-|\alpha|^2}
 \le\frac{M_g(r)^2-M_f(r)^2}{1-\rho^2}
 \le G(r^2)(M_g(r)^2-M_f(r)^2).
\]
The center can be complex and may depend on $r$.
\end{proof}

The following is Wikstr\"om's moment-duality argument~\cite[Lemma 3.1]{Wikstrom},
with the coefficient estimate just proved substituted for its analytic input.
\begin{lemma}\label{lem:mom}
If a probability measure $\mu$ on $[c^2,1]$, with no atom at $1$, satisfies
\[
 \int t^k\,d\mu(t)\ge\frac1{k+1}\quad(k\ge0),\qquad
 \int\beta(t)t^k\,d\mu(t)\le\frac1{k+1}\quad(k\ge1),
\]
then $c$ is admissible in the definition of $c_2$.
\end{lemma}
\begin{proof}
Write $x_k=(|a_k|^2-|b_k|^2)_+$, $y_k=(|b_k|^2-|a_k|^2)_+$,
$X(t)=\sum_{k\ge0}x_kt^k$ and $Y(t)=\sum_{k\ge1}y_kt^k$.
Lemma~\ref{lem:coeff} rearranges to $X(t)\le y_0+\beta(t)Y(t)$.
Integrate and apply the two moment inequalities. Nonnegative-series monotone
convergence gives
$\sum x_k/(k+1)\le y_0+\sum_{k\ge1}y_k/(k+1)$, which is the norm comparison.
Limits give the endpoint convention at $c^2$.
\end{proof}

\section{The rational certificate}
Take $c=54/125$, $R=999/1000$ and
\[
 d\mu(t)=\sum_{j=1}^{12}w_j\,d\delta_{t_j}(t)+\mathbf1_{[R,1]}(t)\,dt.
\]
The entries below are exact terminating decimals, not rounded displays.
They have $t_1=c^2$, $t_{12}=R$, positive weights summing exactly to $R$,
and no atom at $1$.
\begin{center}\small
\begin{tabular}{rll}\toprule
$j$ & $t_j$ & $w_j$\\\midrule
1 & 0.186624 & 0.401004850339305473 \\
2 & 0.587760864611916678 & 0.334589004564615711 \\
3 & 0.8229718555012977 & 0.152356299682819746 \\
4 & 0.924010443899645262 & 0.062924486526793319 \\
5 & 0.965990583618882126 & 0.026544759438047238 \\
6 & 0.984044289625064295 & 0.011714716766935399 \\
7 & 0.992173624211215355 & 0.00540693721877966 \\
8 & 0.995987383331351062 & 0.002581897617100729 \\
9 & 0.997819789755709063 & 0.001239794794409609 \\
10 & 0.998676970469441039 & 0.000541241719003444 \\
11 & 0.998976894402894056 & 0.000064073011221399 \\
12 & 0.999 & 0.000031938320968273 \\
\bottomrule\end{tabular}
\end{center}

Let $g(t)=1/G(t)$. A simple global linear lower bound on the Lebesgue tail
suffices. Define $\Delta=1-R$, $U=-\pi\log R/(2w)$ and
\[
 a=\frac{(\pi/(2w))^2(1-U^2/6)^2}{X},\qquad
 \ell=\frac{2563746486913}{10^{14}}.
\]
Both checkers prove $U^2<6$ and
\begin{equation}\label{eq:ell}
 0<\ell<\frac{2\sqrt{6a}}
 {\sqrt{3a\Delta^2(2+a\Delta^2)}+
  \sqrt{4+6a\Delta^2+3a^2\Delta^4}}.
\end{equation}
Indeed the latter right side is greater than $0.0256374648691356266$.
For $t\in[R,1]$, put $\delta=1-t$ and
$y=\sin^2(\pi\log t/(2w))/X$. Since $-\log t\ge\delta$ and
$\sin v\ge v-v^3/6$, one has $y\ge a\delta^2$. Direct algebra gives
\[
 g(t)=\frac{2\sqrt{3y(2+y)}}{\sqrt{3y(2+y)}+\sqrt{4+6y+3y^2}}.
\]
This expression increases with $y$. Substitute $a\delta^2$ for $y$,
bound its numerator below by $2\sqrt{6a}\delta$ and its denominator above
by the denominator in~\eqref{eq:ell}. Thus $g(t)\ge\ell(1-t)$ on $[R,1]$.
In particular,
\[
 \int_R^1g(t)t^k\,dt\ge\ell\left(
 \frac{1-R^{k+1}}{k+1}-\frac{1-R^{k+2}}{k+2}\right).
\]

The two desired moment inequalities reduce respectively to
\begin{align*}
 A_k&=\sum_jw_jt_j^k-\frac{R^{k+1}}{k+1}\ge0,\\
 B_k&=\int_R^1g(t)t^k\,dt+\frac{R^{k+1}}{k+1}
       -\sum_jw_j\beta(t_j)t_j^k\ge0.
\end{align*}
Here $A_0=0$ by exact mass balance. Both checkers verify $A_k>0$ through
$k=31278$. For all $k\ge31279$, the last atom alone suffices, since
$w_{12}(k+1)\ge R$.

Both checkers verify $B_k>0$ for $1\le k\le23999$ using the displayed
linear integral lower bound and upper enclosures of the atomic $\beta$.
For the rest, set $t_0=0.9998$, $t_1=0.9999$. The exact comparison
\[
 \ell(1-t_1)(t_1-t_0)t_0^{24000}>\sum_jw_jt_j^{24000}
\]
has a certified gap greater than $2.1063\cdot10^{-12}$. Since every
$t_j\le R<t_0$, it persists at every larger exponent. As $\beta\le1$,
it proves the second infinite tail. The verified finite relative margins,
normalized by $R^{k+1}/(k+1)$, exceed $2.0961\cdot10^{-6}$ for $A_k$
and $2.2847\cdot10^{-7}$ for $B_k$. Lemma~\ref{lem:mom} proves
Theorem~\ref{thm:main}.

\section{Independent checks, provenance and limitations}
The supplied \texttt{certificate-jung.json} is exact rational input.
\texttt{verify\_jung\_arb.py} uses Arb at 75 decimal digits.
\texttt{verify\_jung\_rational.py} uses only standard-library integers,
fractions and integer square roots, rounding every operation outwards on a
$2^{200}$ denominator. It implements Machin's formula for $\pi$, the positive
atanh series for $-\log t$, sine with a Lagrange remainder, and exponential
with a geometric remainder after argument reduction. Neither checker imports
the search or the other checker. Seventeen independent arithmetic enclosure
controls against 100-digit Arb pass. Both reject mass-preserving bad moments,
an unsafe linear envelope, a false tail cutoff, a wrong analytic model and
execution with assertions disabled.

The frozen certificate SHA-256 is
\begin{center}\ttfamily\footnotesize
80a056f2cec4cfd933944fa51100bc5a1a9657f1a1de509e3a51401ff282002b.
\end{center}
Its numerical hypotheses, rather than the analytic proof itself, are checked
by the programs. No proof-assistant formalization is claimed. The complete
proof of the substituted analytic estimate is given above for mathematical
review. A separate retained sixteen-atom certificate gives $0.42638465$ using
Wang's original estimate~\cite{Wang2008,Wang} and Wikstr\"om's original moment lemma unchanged.
The exact $0.4263$ public baseline was reproduced before optimization, including
its original published software~\cite{Zenodo}.

Credit goes to Wang and Schuster for the annular-pseudodistance approach,
to the classical Schwarz and hyperbolic Jung inequalities, and to Wikstr\"om
for moment duality, the atomic-plus-Lebesgue construction and its LP/refinement
strategy. The computations widen the atomic support and tail cutoff, then
replace the analytic estimate by the covering-strip/enclosing-disk argument.
Larger unsuccessful numerical trials do not prove a barrier. No optimality
of the measure, analytic estimate or constant is asserted. The analytic proof
and its computational hypotheses remain subject to independent review.

\paragraph{AI use and author review.}
Analysis and proof drafting by AI agents under my direction; I reviewed the
results at summary level. AI agents also prepared the certificate, checking
software, and publication materials. This disclosure does not claim a
line-by-line human verification of the proof or software.

\paragraph{Availability.}
The certificates, checkers, verification receipts, and source are available at
\url{https://github.com/rpinho/korenblum-c68-certificate} under the MIT license.

\begin{thebibliography}{99}
\bibitem{Korenblum} B. Korenblum, \emph{A maximum principle for the Bergman
space}, Publ. Mat. \textbf{35} (1991), no. 2, 479--486.
\bibitem{Hayman} W. K. Hayman, \emph{On a conjecture of Korenblum},
Analysis (Munich) \textbf{19} (1999), no. 2, 195--205.
\bibitem{Wikstrom} F. Wikstr\"om, \emph{Moment duality and an improved lower
bound for Korenblum's constant}, arXiv:2607.17748 (2026),
\url{https://arxiv.org/abs/2607.17748}.
\bibitem{Zenodo} F. Wikstr\"om, \emph{Rigorous certificates for a lower bound on
Korenblum's constant}, Zenodo, doi:10.5281/zenodo.21446369.
\bibitem{Wang2008} C. Wang, \emph{Domination in the Bergman space and
Korenblum's constant}, Integral Equations Operator Theory \textbf{61}
(2008), 423--432, doi:10.1007/s00020-008-1587-4.
\bibitem{Wang} C. Wang, \emph{A lower bound on Korenblum's constant},
Proc. Amer. Math. Soc. \textbf{153} (2025), 1209--1214,
doi:10.1090/proc/17112.
\bibitem{Schuster} A. Schuster, \emph{The maximum principle for the Bergman
space and the M\"obius pseudodistance for the annulus}, Proc. Amer. Math.
Soc. \textbf{134} (2006), 3525--3530; arXiv:math/0407416.
\bibitem{BMRS} S. Baguley, Y. Maus, J. Ruff and G. Skretas,
\emph{Hyperbolic random graphs: clique number and degeneracy with implications
for colouring}, STACS 2025, LIPIcs \textbf{327}, article 13, Theorem 17
(hyperbolic Jung, citing Dekster), doi:10.4230/LIPIcs.STACS.2025.13.
\end{thebibliography}
\end{document}
